\begin{center}
\begin{minipage}{\linewidth}
\captionof{table}{Compared methods and target-side protocols.}
\label{tab:fairness}
\centering
\scriptsize
\renewcommand{\arraystretch}{1.08}
\setlength{\tabcolsep}{1.8pt}
\begin{tabularx}{\linewidth}
{Y Z{0.16\linewidth} Z{0.20\linewidth} Z{0.16\linewidth}
 Z{0.10\linewidth} Z{0.15\linewidth}}
\toprule
Method
& Structure
& Source initialization
& Target update
& Adapted models
& Complexity-limit check \\
\midrule
PT+FT
& Fixed reference architecture
& Pooled source pretraining
& SGD/MSE-50
& 1
& Before output \\

MeDeT-based adaptation
& Fixed reference architecture
& Meta initialization
& SGD/MSE-50
& 1
& Before output \\

Random initialization
& Fixed reference architecture
& Random initialization
& SGD/MSE-50
& 1
& Before output \\

Few-shot NAS
& Top-12 shortlist
& Baseline-specific initialization bank
& Adam/Huber-50
& $\leq 12$
& Before candidate evaluation \\

Zero-shot NAS
& Top-ranked model
& Baseline-specific initialization
& None
& 0
& Before candidate evaluation \\

Zero-shot NAS + fine-tuning
& Top-12 shortlist
& Baseline-specific initialization bank
& Adam/Huber-50
& $\leq 12$
& Before candidate evaluation \\

RCF-DTI
& Source-initialization bank (six architectures)
& Architecture-matched source initializations
& SGD/MSE-50
& 4--7
& Before adaptation and output \\
\bottomrule
\end{tabularx}
\vspace{2pt}
\footnotesize
PT+FT and RCF-DTI use the same source-pretrained reference model and the same
target adaptation protocol.
\end{minipage}
\end{center}
